% =============================================================
%  Kapitel 2 – Gesamtarchitektur
% =============================================================
\section{Gesamtarchitektur des Systems}
\label{sec:gesamtarchitektur}

\subsection{Systemkontext}
\label{sec:systemkontext}

Das Fahrzeug besteht aus drei physisch getrennten Recheneinheiten,
die über definierte Kommunikationskanäle miteinander verbunden sind:

\begin{itemize}
  \item \textbf{Jetson Nano} – Hochleistungs-Recheneinheit für Wahrnehmung und autonome Logik.
  \item \textbf{Arduino Mega Master} – Kommunikations-Gateway zwischen Jetson und Slave.
  \item \textbf{Arduino Mega Slave} – Sicherheitskritische Einheit nahe der Motorsteuerung;
        hier läuft der Emergency Stop.
\end{itemize}

Die Kommunikation erfolgt über:
\begin{itemize}
  \item \textbf{LAN} (Jetson Nano $\leftrightarrow$ Arduino Mega Master)
  \item \textbf{CAN-Bus} (Arduino Mega Master $\leftrightarrow$ Arduino Mega Slave)
\end{itemize}

Die Platzierung des Emergency Stop auf dem Slave minimiert Kommunikationsverzögerungen
zur Motorsteuerung (\emph{DriveControlUnit}) und stellt eine deterministische
Reaktionszeit sicher.

\begin{figure}[H]
  \centering
  \includegraphics[width=0.95\textwidth]{01_Context}
  \caption{Kontextdiagramm – Systemgrenzen und externe Akteure}
  \label{fig:context}
\end{figure}

\subsection{Logische Systemstruktur (Component Overview)}
\label{sec:component_overview}

Die logische Architektur zeigt die funktionale Zerlegung des Systems in Subsysteme.
Abbildung~\ref{fig:component_overview} gibt einen Überblick.

\begin{figure}[H]
  \centering
  \includegraphics[width=0.95\textwidth]{02_Component_Overview}
  \caption{Komponentenübersicht – logische Systemstruktur}
  \label{fig:component_overview}
\end{figure}

Wesentliche Eigenschaften der Architektur:
\begin{itemize}
  \item Der \textbf{TTC-Observer} verarbeitet Sensordaten zur Bestimmung der Fahrzone
        (Grün/Gelb) und gehört ausschließlich zur \emph{Normal Operation}.
  \item Der \textbf{Emergency Stop} verarbeitet ausschließlich CAN-basierte Kommandos
        und ist von der Autonomie-Logik entkoppelt.
  \item Sensorik wird über das \texttt{EnvironmentControl}-Interface in CAN-Nachrichten
        übersetzt (Sensor-Fusion als Blackbox).
  \item Das \textbf{Logging}-Modul ist als separates Querschnittsmodul modelliert;
        es wird von Emergency Stop bei jeder Zustandsänderung aufgerufen.
  \item Eine direkte Abhängigkeit zwischen TTC-Observer und Emergency Stop
        \textbf{existiert nicht} – dies sichert die Unabhängigkeit der Sicherheitsfunktion.
\end{itemize}

\subsection{Subsystem-Verantwortlichkeiten und Klassendiagramme}
\label{sec:klassendiagramme}

Die folgenden Klassendiagramme zeigen Attribute und öffentliche Operationen
pro Hardware-Subsystem.

\subsubsection{Jetson Nano – Klassen}

\begin{figure}[H]
  \centering
  \includegraphics[width=0.95\textwidth]{02_Jetson_ClassDiagram}
  \caption{Klassendiagramm Jetson Nano: Autonomie-/Wahrnehmungs-Subsystem}
  \label{fig:jetson_classes}
\end{figure}

Hauptverantwortlichkeiten des Jetson-Subsystems:
\begin{description}
  \item[\texttt{AutonomyController}]
    Koordiniert autonomes Fahrverhalten; empfängt Zonenstatus vom TTC-Observer.
  \item[\texttt{TTCObserver}]
    Berechnet Time-To-Collision aus Sensordaten; sendet ZoneUpdate-Kommandos
    (Normal/Warnung). Zuständig für Gelb- und Grün-Zonen.
  \item[\texttt{SensorFusionAdapter}]
    Abstrahiert die Sensor-Fusion (Kamera + LiDAR); liefert normierte
    Distanz-/Objektdaten. Intern als Blackbox behandelt.
\end{description}

\subsubsection{Arduino Mega Master – Klassen}

\begin{figure}[H]
  \centering
  \includegraphics[width=0.95\textwidth]{02_Master_ClassDiagram}
  \caption{Klassendiagramm Arduino Mega Master: Kommunikations-Gateway}
  \label{fig:master_classes}
\end{figure}

Hauptverantwortlichkeiten des Master-Subsystems:
\begin{description}
  \item[\texttt{MasterLanLink}]
    Empfängt Befehle vom Jetson über LAN und leitet sie auf den CAN-Bus weiter.
  \item[\texttt{MasterCanGateway}]
    Kapselt CAN-Kommunikation; wandelt LAN-Nachrichten in CAN-Pakete um.
    Nutzt intern die Klasse \texttt{CanMessage} (siehe unten).
  \item[\texttt{HmiService}]
    Verarbeitet Signale des Pilzknopfs (Emergency-Button) und des Fronttasters;
    erzeugt entsprechende CAN-Ereignisse.
\end{description}

\subsubsection{Datenklasse: CanMessage}

Für die CAN-Kommunikation wird gemäß OOP-Anforderung eine eigene
Datenklasse \texttt{CanMessage} eingesetzt, die auf Master- und
Slave-Seite identisch verwendet wird:

\begin{center}
\renewcommand{\arraystretch}{1.3}
\begin{tabular}{>{{\ttfamily}}p{3.5cm} >{{\ttfamily}}p{2.8cm} p{7cm}}
  \toprule
  \normalfont\textbf{Attribut} & \normalfont\textbf{Typ} & \normalfont\textbf{Bedeutung} \\
  \midrule
  messageId   & uint16\_t  & CAN-Identifier; kodiert den Nachrichtentyp
                             (z.\,B.\ ZoneUpdate, EmergencyEvent). \\
  payload     & uint8\_t[8] & Nutzdaten (bis zu 8\,Byte, CAN-Standard). \\
  payloadLen  & uint8\_t   & Anzahl gültiger Payload-Bytes. \\
  timestamp   & uint32\_t  & Zeitstempel in Millisekunden seit Systemstart. \\
  \bottomrule
\end{tabular}
\end{center}

\texttt{CanMessage} ist eine reine Datenklasse (\emph{Plain Old Data})
ohne eigene Logik; sie trennt die Transportschicht von der
Anwendungsschicht und ermöglicht ein einheitliches Logging.

\subsubsection{Arduino Mega Slave – Klassen}

\begin{figure}[H]
  \centering
  \includegraphics[width=0.95\textwidth]{02_Slave_ClassDiagram}
  \caption{Klassendiagramm Arduino Mega Slave: Sicherheits- und Antriebssubsystem}
  \label{fig:slave_classes}
\end{figure}

Hauptverantwortlichkeiten des Slave-Subsystems:
\begin{description}
  \item[\texttt{EnvironmentControl}]
    Schnittstelle zur Empfangsseite des CAN-Busses; übersetzt eingehende Pakete
    in interne Ereignisse für den Emergency Stop.
  \item[\texttt{EmergencyStop}]
    Zentrale Sicherheitsklasse (Detailed Design in Abschnitt~\ref{sec:detaileddesign}).
    Steuert den Zustandsautomaten und die DriveControlUnit.
  \item[\texttt{DriveControlUnit}]
    Abstrahiert die Motorsteuerung; empfängt ausschließlich Stop-/Freigabe-Befehle
    vom Emergency Stop.
  \item[\texttt{Logger}]
    Querschnittsmodul für Error-Logging; unterstützt Level (\texttt{ERROR},
    \texttt{WARNING}, \texttt{INFO}) mit Zeitstempel.
\end{description}
